\subsection{Procesoriaus modelis}
\label{sec:procesoriaus-modelis}

Procesoriaus modelis projektuotas taip, kad tą patį vykdymo branduolį būtų galima naudoti kelioms RISC-V konfigūracijoms. Kompiliavimo metu parenkamas žodžio dydis ir įjungiami reikalingi ISA plėtiniai: bazinis RV32 arba RV64 instrukcijų rinkinys, privilegijuotas vykdymas, M ir A plėtiniai bei eksperimentinis QPU plėtinys.

Projektuojant procesoriaus būseną buvo išskirtos kelios grupės:

\begin{itemize}
    \item \textbf{registrai} -- 32 bendrosios paskirties registrai (nuo \texttt{x0} iki \texttt{x31}) ir programos skaitiklis \texttt{pc}. Registras \texttt{x0} visada išlaikomas lygus nuliui;
    \item \textbf{privilegijuota būsena} -- dabartinis privilegijų režimas, CSR būsena; 
    \item \textbf{atominės operacijos} -- standartinio A plėtinio LR/SC rezervacija;
    \item \textbf{QPU būsena} -- procesoriuje saugomas aktyvus QPU konteksto identifikatorius, o kvantinių registrų būsena laikoma atskiroje QPU posistemėje.
\end{itemize}

\subsubsection{Instrukcijų dekodavimas}

Emuliatorius palaikys tik 32 bitų instrukcijas -- suspaustosios 16 bitų instrukcijos nepalaikomos. Dekoduojant instrukciją \texttt{instrukcijos\_zodis}, ji yra išskaidoma į bitų laukus, pavaizduotus \ref{table:funct-laukai} lentelėje. Pagrindinis laukas yra \texttt{opcode}, o papildomam atskyrimui naudojami \texttt{funct3}, \texttt{funct5}, \texttt{funct6} ir \texttt{funct7} laukai.

\begin{table}[h]
    \centering
    \small
    \begin{tabular}{|l|l|}
        \hline
        \textbf{Laukas} & \textbf{Priskiriama reikšmė} \\
        \hline
        \texttt{opcode} & \texttt{instrukcijos\_zodis[6:0]} \\
        \hline
        \texttt{funct3} & \texttt{instrukcijos\_zodis[14:12]} \\
        \hline
        \texttt{funct5} & \texttt{instrukcijos\_zodis[31:27]} \\
        \hline
        \texttt{funct6} & \texttt{instrukcijos\_zodis[31:26]} \\
        \hline
        \texttt{funct7} & \texttt{instrukcijos\_zodis[31:25]} \\
        \hline
    \end{tabular}
    \caption{Instrukcijos žodžio laukai, naudojami instrukcijų atpažinimui.}
    \label{table:funct-laukai}
\end{table}

Visų pirma, surandamos visos potencialios instrukcijos, kurių opkodas sutampa. Tuomet kiekviena instrukcija gali turėti papildomų kaukių, kurios privalo sutapti. Palyginimo strategija nusako, kurie iš \texttt{instrukcijos\_zodis} išskirti laukai turi sutapti su instrukcijos apraše nurodytomis reikšmėmis. Pavyzdžiui, jei instrukcija atpažįstama pagal \texttt{funct3} ir \texttt{funct7}, tada lyginamas ne visas instrukcijos žodis, o tik \texttt{opcode}, \texttt{funct3} ir \texttt{funct7} reikšmės. Kiekviena instrukcija turi priskirtą palyginimo strategiją, kurios išvardintos \ref{table:palyginimo-strategijos} lentelėje.

\begin{table}[h]
    \centering
    \small
    \begin{tabular}{|l|c|c|c|c|c|c|c|c|}
        \hline
        \diagbox[width=10em]{\textbf{Kaukė}}{\textbf{Instrukcijos}\\\textbf{bitai}}
        & \textbf{31:27}
        & \textbf{26}
        & \textbf{25}
        & \textbf{24:20}
        & \textbf{19:15}
        & \textbf{14:12}
        & \textbf{11:7}
        & \textbf{6:0} \\
        \hline

        \texttt{opcode}
        & -- & -- & -- & -- & -- & -- & -- & \texttt{+} \\
        \hline

        \texttt{funct3 + funct7}
        & \texttt{+}
        & \texttt{+}
        & \texttt{+}
        & -- & -- & + & -- & \texttt{+} \\
        \hline

        \texttt{funct3 + funct6}
        & \texttt{+}
        & \texttt{+}
        & -- & -- & -- & + & -- & \texttt{+} \\
        \hline

        \texttt{funct3 + funct5}
        & \texttt{+}
        & -- & -- & -- & -- & + & -- & \texttt{+} \\
        \hline

        Tiesioginė
        & \multicolumn{8}{c|}{\texttt{+}} \\
        \hline
    \end{tabular}
    \caption{Instrukcijų atpažinimo kaukės. Pliusas reiškia, kad atitinkami bitai yra lyginami, brūkšnys -- nelyginami.}
    \label{table:palyginimo-strategijos}
\end{table}
